% 本文件由 tools/build_book.py 从 dist/reports/assertions.json 生成，请勿手改。
\begin{longtable}{@{}p{0.06\textwidth}p{0.32\textwidth}p{0.22\textwidth}p{0.16\textwidth}@{}}
\caption{章节 5.4 数值断言（生成时间 2026-10-18T18:30:41-0400）}\\
\toprule
类别 & 断言 & 理论值 & 实测值 \\
\midrule
\endfirsthead
\toprule
类别 & 断言 & 理论值 & 实测值 \\
\midrule
\endhead
\bottomrule
\endlastfoot

A & 端点组合帧数与 JSON 里真正的非有限常量个数 & 非有限常量个数 =\allowbreak{} 0(\allowbreak{}字符串里出现 NaN 三个字母不算)，帧数 >=\allowbreak{} 60 & 66 帧，计数 0(\allowbreak{}point\_fill: 17 个数组、\allowbreak{}17 个 meta 字段) \checkmark \\
A & 对数曲线范围、地板说明、以及贴地板处掩码确实为 0 & 全部曲线 >=\allowbreak{} -17 且 max > -16；discontinuous 掩码全 0、\allowbreak{}offset\_box 掩码全 1 & point\_fill.error$\in$[-5.5,\allowbreak{}-3.5]；point\_fill.error\_mc$\in$[-2.8,\allowbreak{}-1.8]；point\_fill.error\_sobol$\in$[-5.5,\allowbreak{}-3.5]；point\_fill.error\_ref$\in$[-2.7,\allowbreak{}-1.8]；point\_fill.band\_lo$\in$[-5.6,\allowbreak{}-3.6]；point\_fill.band\_hi$\in$[-5.4,\allowbreak{}-3.4]；output\_hist.error$\in$[-2.9,\allowbreak{}-1.0]；output\_hist.bins\_error$\in$[-1.3,\allowbreak{}-1.3]；pce\_spectrum.error$\in$[-11.5,\allowbreak{}-0.6]；pce\_spectrum.error\_train$\in$[-11.6,\allowbreak{}-0.7]；pce\_spectrum.coeff\_rms$\in$[-3.5,\allowbreak{}0.2]；discontinuous 掩码 [0.]，offset\_box 掩码 [1.] \checkmark \\
A & 三种情形下的 band\_\allowbreak{}valid / band\_\allowbreak{}available / 文字说明 & 关随机化与 grid：全 0 且 band\_available=\allowbreak{}False；开随机化：全 1 且 True & 关：valid=\allowbreak{}[0.]，available=\allowbreak{}False；grid：valid=\allowbreak{}[0.]，available=\allowbreak{}False；开：valid=\allowbreak{}[1.]，available=\allowbreak{}True；说明含「不存在」=\allowbreak{}True \checkmark \\
A & 贴地板时的 slope\_\allowbreak{}current 与说明文字 & slope\_current =\allowbreak{} None；note 里出现「精确为 0」；偏移边界上斜率 $\in$ [-1.2,\allowbreak{}-0.6] & slope\_current=\allowbreak{}None，note 片段=\allowbreak{}'sobol 的误差精确为 0(\allowbreak{}点集与间断/\allowbreak{}边界对齐，见 error\_valid)'，offset\_box 的 slope\_current=\allowbreak{}-0.8698731004959773 \checkmark \\
A & rho=0 与 rho=0.5 的字段对照(双线性模型) & rho=\allowbreak{}0：三个字段为 None、\allowbreak{}corr\_valid 全 0；rho=\allowbreak{}0.5：sum\_between > 1.4 且 sum\_indep $\in$ [0.95,\allowbreak{}1.05] & rho=\allowbreak{}0：corr\_valid=\allowbreak{}[0.]，overlap=\allowbreak{}None，sum\_between=\allowbreak{}None；rho=\allowbreak{}0.5：corr\_valid=\allowbreak{}[1.]，sum\_between=\allowbreak{}1.4942，sum\_indep=\allowbreak{}0.9890，rho\_measured=\allowbreak{}0.4990，corr\_note 带数字=\allowbreak{}True \checkmark \\
A & 均匀输入方体上的模型拖 rho & corr\_applicable=\allowbreak{}False、\allowbreak{}corr\_active=\allowbreak{}False，且文字说明 rho 无效 & corr\_applicable=\allowbreak{}False，corr\_active=\allowbreak{}False，说明含「无效」=\allowbreak{}True \checkmark \\
A & 独立/相关两种情形下的和与判据布尔 & 独立：sum S\_i <=\allowbreak{} 1 <=\allowbreak{} sum S\_Ti 且布尔为真；相关：sum\_between > 1(\allowbreak{}判据在该口径下失效) & 独立 g-函数：sum S\_i=\allowbreak{}0.9348，sum S\_Ti=\allowbreak{}1.1077，lower=\allowbreak{}True，upper=\allowbreak{}True，解析=\allowbreak{}0.9290；相关双线性：sum\_between=\allowbreak{}1.4942(\allowbreak{}>1) \checkmark \\
A & 过拟合区的三个字段 + 训练/留出误差 & underdetermined=\allowbreak{}True；under\_flag =\allowbreak{} (\allowbreak{}n\_basis > 8) 逐点成立；训练 <=\allowbreak{} 1e-14 且留出 > 0.1 & underdetermined=\allowbreak{}True，欠定阶数 5，训练 4.49e-16，留出 4.855e-01，外推 1.392e+03(\allowbreak{}放大 2.87e+03 倍) \checkmark \\
A & poly4 的 exact\_\allowbreak{}at\_\allowbreak{}order / 掩码 / 当前阶误差 & exact\_at\_order =\allowbreak{} 4；p<4 时 log10(\allowbreak{}err) >=\allowbreak{} -1.1；当前阶误差 <=\allowbreak{} 1e-12；掩码全 1 & exact\_at\_order=\allowbreak{}4，log10 误差曲线 [ -0.36  -0.89  -1.05 -14.6  -14.75 -14.38 -14.4  -14.45 -13.73 -13.81
 -13.41 -13.78]，掩码取值 [1.]，当前留出误差 1.651e-14 \checkmark \\
A & 三种输入分布下 ref\_\allowbreak{}is\_\allowbreak{}exact / KS 对照 & 正态+线性：exact=\allowbreak{}True 且两个 KS 相差 <=\allowbreak{} 2e-3(\allowbreak{}精确参考就是正态)；均匀+线性：精确三角密度，ks\_to\_normal >=\allowbreak{} 3 * ks\_to\_ref；指数+线性：没有精确参考，exact=\allowbreak{}False 且 ref\_kind 里写明是 CLT 近似 & 正态：exact=\allowbreak{}True，ks\_ref=\allowbreak{}0.0064，ks\_normal=\allowbreak{}0.0054；均匀：exact=\allowbreak{}True，ks\_ref=\allowbreak{}0.0029，ks\_normal=\allowbreak{}0.0190；指数：exact=\allowbreak{}False，ref\_kind=\allowbreak{}'CLT 正态近似(\allowbreak{}同均值同方差的正态密度，不是真' \checkmark \\
A & 四处数组与核心调用的逐位/对数一致性 & 点集与指数逐位相等(\allowbreak{}差 =\allowbreak{} 0)；计数和 =\allowbreak{} N；误差对数差 <=\allowbreak{} 1e-12 & 点集差 0.0e+00，指数差 0.0e+00，计数和 20000，误差 log10 差 0.000e+00(\allowbreak{}帧的曲线来自同一个 pce.error\_vs\_order 调用) \checkmark \\
A & norm\_\allowbreak{}ppf 与 verify.stats 的最大偏差 + cdf 往返 + 端点行为 & 偏差 =\allowbreak{} 0(\allowbreak{}同一有理逼近)、\allowbreak{}往返 <=\allowbreak{} 1e-8、\allowbreak{}端点为 $\pm$inf & 偏差 0.000e+00，往返 2.465e-10，ppf(\allowbreak{}0)=\allowbreak{}-inf，ppf(\allowbreak{}1)=\allowbreak{}inf \checkmark \\
A & 四种分布的一阶/二阶矩与支撑集 & |mean| <=\allowbreak{} 0.02、\allowbreak{}|std-1| <=\allowbreak{} 0.02 & uniform: mean=\allowbreak{}-0.0009 std=\allowbreak{}1.0006 range=\allowbreak{}[-1.73,\allowbreak{}1.73]；normal: mean=\allowbreak{}+0.0000 std=\allowbreak{}1.0011 range=\allowbreak{}[-4.49,\allowbreak{}4.43]；beta: mean=\allowbreak{}-0.0023 std=\allowbreak{}0.9992 range=\allowbreak{}[-1.79,\allowbreak{}4.08]；exponential: mean=\allowbreak{}-0.0010 std=\allowbreak{}0.9979 range=\allowbreak{}[-1.00,\allowbreak{}12.74]；beta 偏度 0.593(\allowbreak{}Beta(\allowbreak{}2,\allowbreak{}5) 解析偏度 0.596) \checkmark \\
A & 均匀和的经验 L1 距离 / 精确密度积分 / 指数和的返回值 & L1 < 0.02、\allowbreak{}积分 =\allowbreak{} 1(\allowbreak{}误差 <=\allowbreak{} 1e-4 是 60 格黎曼和的求积误差)、\allowbreak{}指数返回 None & L1=\allowbreak{}0.01402，积分=\allowbreak{}0.999981，指数分布的返回值=\allowbreak{}None \checkmark \\
B & 直方图密度积分 / 计入样本数 / bin 误差曲线形状 & 积分 =\allowbreak{} 1(\allowbreak{}1e-12)、\allowbreak{}计数和 =\allowbreak{} N、\allowbreak{}两端误差 > 内部极小 & 积分=\allowbreak{}1.000000000000，计数 50000/\allowbreak{}50000，极小在 B=\allowbreak{}16(\allowbreak{}误差 0.0029)，B=\allowbreak{}8 $\to$ 0.0052，B=\allowbreak{}96 $\to$ 0.0062 \checkmark \\
C & MC 的双对数斜率(框架 fit\_\allowbreak{}order 与本模块 loglog\_\allowbreak{}slope) & 0.5 $\pm$ 0.08，且两条路径一致 & fit\_order 0.4731(\allowbreak{}R$^2$=\allowbreak{}0.987073，n\_used=\allowbreak{}6)，loglog\_slope 0.4731，误差序列 [0.01365 0.01112 0.00848 0.00572 0.00416 0.00268 0.00194 0.00134] \checkmark \\
C & N=100 时两种输出的 95\% 覆盖率与其差 & 正态输出 $\in$ [0.90,\allowbreak{}0.99]；重尾输出低于它(\allowbreak{}差距在二项标准误 0.013 的 2 倍量级) & 正态 0.9533(\allowbreak{}命中 286/\allowbreak{}300)vs 重尾 0.9267(\allowbreak{}命中 278/\allowbreak{}300)，差 +0.0267(\allowbreak{}se$\approx$0.0122)，半宽 0.2781 vs 0.2856 \checkmark \\
C & N=1000 时同一个重尾输出的覆盖率(差距是否消失) & 两者都回到 0.95 附近(\allowbreak{}|覆盖率 - 0.95| <=\allowbreak{} 3 se) & 正态 0.9467，重尾 0.9567，|重尾-0.95|=\allowbreak{}0.0067(\allowbreak{}3se=\allowbreak{}0.0353) \checkmark \\
B & 精确参考下的 KS 序列 / CLT 参考下的 KS 序列 & 精确参考：KS(\allowbreak{}32000) < KS(\allowbreak{}2000)；CLT 参考：KS 不随 N 消失(\allowbreak{}地板 > 5e-3) & 精确 [0.01492 0.00722 0.00422](\allowbreak{}N=\allowbreak{}(\allowbreak{}2000,\allowbreak{} 8000,\allowbreak{} 32000))，CLT [0.0265  0.02037 0.01766] \checkmark \\
A & Sobol 网三种置乱下的每格点数极值 vs MC & Sobol：min=\allowbreak{}max=\allowbreak{}1(\allowbreak{}所有 m1+m2=\allowbreak{}m 的切法)；MC：min=\allowbreak{}0 且 max >=\allowbreak{} 2 & m=\allowbreak{}4/\allowbreak{}none: [1,\allowbreak{}1]；m=\allowbreak{}4/\allowbreak{}shift: [1,\allowbreak{}1]；m=\allowbreak{}4/\allowbreak{}linear: [1,\allowbreak{}1]；m=\allowbreak{}6/\allowbreak{}none: [1,\allowbreak{}1]；m=\allowbreak{}6/\allowbreak{}shift: [1,\allowbreak{}1]；m=\allowbreak{}6/\allowbreak{}linear: [1,\allowbreak{}1]；m=\allowbreak{}8/\allowbreak{}none: [1,\allowbreak{}1]；m=\allowbreak{}8/\allowbreak{}shift: [1,\allowbreak{}1]；m=\allowbreak{}8/\allowbreak{}linear: [1,\allowbreak{}1]；MC 256 点在 16x16 格上 [0,\allowbreak{}4](\allowbreak{}空格 86 个) \checkmark \\
A & HK 不等式被违反的次数 / 松紧度的分布 & 违反 0 次；松紧度 min < 1e-2 且 max > 0.5 & 违反 0/\allowbreak{}20 次；松紧度 min=\allowbreak{}0.0000，max=\allowbreak{}0.9685，中位数=\allowbreak{}0.0288(\allowbreak{}跨 300.0 个数量级) \checkmark \\
A & 变差的频率标度 V(k)/(V(1)k\textasciicircum{}2)(两种网格分辨率) & =\allowbreak{} 1；400 格上偏差 <=\allowbreak{} 2e-3(\allowbreak{}差分离散化)，1600 格上再小 4 倍以上 & V(\allowbreak{}2)/\allowbreak{}(\allowbreak{}V(\allowbreak{}1)*2\textasciicircum{}2)=\allowbreak{}0.999814466；V(\allowbreak{}4)/\allowbreak{}(\allowbreak{}V(\allowbreak{}1)*4\textasciicircum{}2)=\allowbreak{}0.999072560；V(\allowbreak{}8)/\allowbreak{}(\allowbreak{}V(\allowbreak{}1)*8\textasciicircum{}2)=\allowbreak{}1.000067037；400 格最大偏差 9.27e-04，1600 格最大偏差 2.43e-04 \checkmark \\
A & D*(37) 的 MC/Sobol 序列与 D*(64) 的地板 & 每个 N 上 D*(\allowbreak{}Sobol) < D*(\allowbreak{}MC)/\allowbreak{}4，且优势随 N 增大；D*(\allowbreak{}64) =\allowbreak{} 0(\allowbreak{}N=\allowbreak{}4096 时，锚网格分辨率地板) & 37 格：MC [0.06153 0.03517 0.01768]，Sobol [0.01152 0.00275 0.00087]；64 格 Sobol [0.00562 0.00073 0.     ] \checkmark \\
B & 确定性复制的标准差 / 真实误差 / 覆盖率 / 随机化后的误差棒 & 标准差 =\allowbreak{} 0、\allowbreak{}误差 > 1e-4、\allowbreak{}覆盖率 =\allowbreak{} 0、\allowbreak{}随机化后 stderr > 0 且 bar\_honesty*sqrt(\allowbreak{}R) $\approx$ 1 & 复制标准差 0.0e+00，误差 2.866e-04，确定性覆盖率 0.00，随机化 stderr\_mean 6.997e-07，bar\_honesty*sqrt(\allowbreak{}R)=\allowbreak{}0.986 \checkmark \\
B & 批量均值与真实标准误的比值(三种方法) & MC $\in$ [0.75,\allowbreak{}1.35]、\allowbreak{}Sobol $\in$ [0.5,\allowbreak{}1.5]、\allowbreak{}格点 $\in$ [0.5,\allowbreak{}1.5] & mc: 批量均值 2.773e-04 /\allowbreak{} 真实标准误 2.426e-04 =\allowbreak{} 1.143；sobol: 批量均值 2.390e-07 /\allowbreak{} 真实标准误 2.952e-07 =\allowbreak{} 0.810；lattice: 批量均值 2.941e-05 /\allowbreak{} 真实标准误 3.084e-05 =\allowbreak{} 0.954 \checkmark \\
C & 小窗口(帧所用)与大窗口上的 MC / Sobol 斜率 & 小窗口：MC $\in$ [-0.58,\allowbreak{}-0.42] 且 Sobol 至少快 0.30；大窗口：Sobol 更陡(\allowbreak{}斜率随窗口变) & 小窗口 [64,\allowbreak{}4096]：MC -0.4969，Sobol -1.2185(\allowbreak{}差 +0.7216)；大窗口 [1024,\allowbreak{}65536]：MC -0.5151，Sobol -1.9034 \checkmark \\
C & 间断边界在 2\textasciicircum{}-1(与网对齐)vs 在 0.3(不对齐) & 对齐：14 个 (\allowbreak{}N,\allowbreak{} 随机化) 组合的误差全为 0；不对齐：斜率 $\approx$ -0.85 且误差非零 & 对齐误差 [0. 0. 0. 0. 0. 0. 0. 0. 0. 0. 0. 0. 0. 0.](\allowbreak{}max=\allowbreak{}0e+00)；不对齐斜率 -0.898，误差 [0.01 0.01 0.   0.   0.   0.   0.  ] \checkmark \\
A & 生成矩阵累加 vs sobol\_\allowbreak{}points 的集合差 / 置乱后的 1 维分层 & 集合差 =\allowbreak{} 0；置乱后 256 点落在 256 个 1 维格子里，每格恰好 1 个 & 集合差 0.000e+00；置乱后 1 维 256 格(\allowbreak{}n=\allowbreak{}256)的计数极值 [1,\allowbreak{}1] \checkmark \\
A & Ishigami 解析值 vs 闭式公式 & 相对差 <=\allowbreak{} 1e-12；sum S\_i =\allowbreak{} 0.7563 < 1 < sum S\_Ti =\allowbreak{} 1.2437 & 最大相对差 2.776e-17；S=\allowbreak{}[0.3139 0.4424 0.    ]，S\_T=\allowbreak{}[0.5576 0.4424 0.2437]，sum S\_i=\allowbreak{}0.7563，sum S\_Ti=\allowbreak{}1.2437，V=\allowbreak{}13.8446 \checkmark \\
A & g-函数解析值 vs 闭式 + 错误公式的对照 & 正确公式相对差 <=\allowbreak{} 1e-12 且 sum S\_i < 1；错误公式 sum S\_i =\allowbreak{} 1.0716 > 1 & 最大相对差 0.000e+00；正确 sum S\_i=\allowbreak{}0.9290，sum S\_Ti=\allowbreak{}1.0716；错误公式 sum S\_i=\allowbreak{}1.0716 \checkmark \\
A & Gram 矩阵偏离单位阵 / 节点与权重对拍 / Hermite 三点权重 & Gram <=\allowbreak{} 1e-9；节点 <=\allowbreak{} 5e-15；权重比恒为 2(\allowbreak{}<=\allowbreak{} 1e-15)；Hermite 三点 =\allowbreak{} (\allowbreak{}1/\allowbreak{}6,\allowbreak{}2/\allowbreak{}3,\allowbreak{}1/\allowbreak{}6) & Gram 最大偏离 4.944e-11；GL 节点 1.887e-15；GL 权重差 6.731e-16；GH 三点权重 [0.166667 0.666667 0.166667] \checkmark \\
A & poly4 在 p=2,3,4,6,12 的投影误差 与 |x| 在 p=12 的误差 & poly4：p<4 时 > 1e-2、\allowbreak{}p>=\allowbreak{}4 时 <=\allowbreak{} 1e-12；|x| 在 p=\allowbreak{}12 时 > 1e-2 & poly4 [1.292e-01 9.158e-02 2.457e-15 4.020e-15 1.627e-14]；|x| p=\allowbreak{}12 $\to$ 0.0319 \checkmark \\
A & PCE 解析指数 vs pick-freeze 实测(双线性，N=4096) & PCE 精确给出 (\allowbreak{}0.5,\allowbreak{}0.5)(\allowbreak{}相对 1e-12)；pick-freeze 在 $\pm$0.05 内 & PCE S=\allowbreak{}[0.5 0.5]，S\_T=\allowbreak{}[0.5 0.5]；pick-freeze S=\allowbreak{}[0.5046 0.5064](\allowbreak{}半宽 [0.0315 0.0306])，V=\allowbreak{}2.0016 \checkmark \\
B & 独立输入：g-函数的 sum S\_\allowbreak{}i / sum S\_\allowbreak{}Ti 与解析值、区间 & sum S\_i - 半宽 <=\allowbreak{} 1(\allowbreak{}不显著超过 1)、\allowbreak{}sum S\_Ti >=\allowbreak{} 1(\allowbreak{}区间内)，且与解析值差 <=\allowbreak{} 0.08 & sum S\_i=\allowbreak{}0.9348$\pm$0.0820(\allowbreak{}解析 0.9290)，sum S\_Ti=\allowbreak{}1.1077$\pm$0.0358(\allowbreak{}解析 1.0716) \checkmark \\
B & 相关输入(rho=0.5)：条件方差口径 / 乘积测度口径 & 条件方差口径 > 1(\allowbreak{}实测 1.4942，超出部分 0.4942)；乘积测度口径 $\approx$ 1(\allowbreak{}实测 0.989) & 条件方差口径 1.4942(\allowbreak{}超出 +0.4942)，乘积测度口径 0.9890，两者差 +0.5052 \checkmark \\
B & sum S\_\allowbreak{}i(N) 的漂移幅度与解析值对照 & 漂移 >=\allowbreak{} 0.15；N=\allowbreak{}4096 处与解析值差 <=\allowbreak{} 0.08 & sum S\_i(\allowbreak{}N)=\allowbreak{}[1.1587 0.8416 1.1657 1.0034 0.8684 0.932  0.9488](\allowbreak{}N=\allowbreak{}[64,\allowbreak{} 128,\allowbreak{} 256,\allowbreak{} 512,\allowbreak{} 1024,\allowbreak{} 2048,\allowbreak{} 4096])，漂移 0.3242，解析 0.9290，末端差 0.0198 \checkmark \\
C & 两种被积函数的半对数衰减率 & exp：斜率 <=\allowbreak{} -0.9 且 R$^2$ >=\allowbreak{} 0.98；|x|：斜率 >=\allowbreak{} -0.2(\allowbreak{}远慢于几何) & exp: 半对数斜率 -1.0005(\allowbreak{}R$^2$=\allowbreak{}0.9919)；abs: 半对数斜率 -0.1080(\allowbreak{}R$^2$=\allowbreak{}0.8596)；exp 从 p=\allowbreak{}1 到 p=\allowbreak{}12 掉了 10.9 个数量级，|x| 只掉了 1.4 个数量级 \checkmark \\
C & |x| 的代数衰减指数(三种子样本) & 三者都落在 [0.9,\allowbreak{}1.5] 且 R$^2$ >=\allowbreak{} 0.97：与渐近的 p\textasciicircum{}-1(\allowbreak{}c\_k \textasciitilde{} k\textasciicircum{}(\allowbreak{}-3/\allowbreak{}2))一致，但 p<=\allowbreak{}16 上测不出干净的 -1(\allowbreak{}奇偶台阶) & 全部阶 alpha=\allowbreak{}1.2439(\allowbreak{}R$^2$=\allowbreak{}0.9811)；奇数阶 alpha=\allowbreak{}1.3061(\allowbreak{}R$^2$=\allowbreak{}0.9996)；偶数阶 alpha=\allowbreak{}1.0940(\allowbreak{}R$^2$=\allowbreak{}0.9972)；误差序列 [0.5099 0.1282 0.1282 0.0638 0.0646 0.0405 0.0405 0.0296 0.0297 0.0217
 0.022  0.0182 0.0188 0.0149 0.0149 0.0144] \checkmark \\
\end{longtable}
